\paragraph{\trMutation{}。}\label{sub:mutation}
对一个给定的\trExceptionalCollection{} $\mathbf{Y}=(Y_1,\ldots,Y_m)$，可以施加
下列变换来修改它：
\begin{itemize}
\item 把每个 $Y_k$ 换成一个在 $H^0(\scr{A})$ 中与它\trIsomorphic{}的%
\trObject{} $\tilde{Y}_k$，并构成所得的序列 $\widetilde{\mathbf{Y}}$；
\item 给定 $(\sigma_1,\ldots,\sigma_m)\in\ZZ^m$，取
$\mathscr{S}^{\sigma}\mathbf{Y}=(S^{\sigma_1}Y_1,\ldots,S^{\sigma_m}Y_m)$，
其中 $S$ 是 $\scr{A}$ 上的\trShiftFunctor；
\item 最后，第 $k$ 个\textbf{初等\trMutation}（elementary mutation）
（$1\le k<m$）是运算 $\mathbf{Y}\mapsto\mathbf{U}=\mathscr{L}_k\mathbf{Y}$，
其中当 $j\neq k,k+1$ 时 $U_j=Y_j$，而剩下的两个\trObject{}则为
\begin{equation}\label{eq:elementary-mutation}
U_k=T_{Y_k}Y_{k+1},\qquad U_{k+1}=Y_k.
\end{equation}
\end{itemize}

目前姑且在\trExceptionalCollection{}的\trIsomorphism{}类这一层面上工作，这就是
说，我们把第一种变换视为平凡的。这时，第二、第三种变换都是\trInvertible{}的。
后者的逆是 $\mathbf{Y}\mapsto\mathbf{V}=\mathscr{L}_k^{-1}\mathbf{Y}$，与前一样，
当 $j\neq k,k+1$ 时 $V_j=Y_j$，而
\begin{equation*}
V_k=Y_{k+1},\qquad V_{k+1}=T^{\vee}_{Y_{k+1}}Y_k.
\end{equation*}

其次是\trBraid{}关系：
\begin{equation*}
\begin{gathered}
\mathscr{L}_j\mathscr{L}_k\mathbf{Y}\cong\mathscr{L}_k\mathscr{L}_j\mathbf{Y}
\quad\text{当 }|j-k|\ge 2\text{ 时，}\\
\mathscr{L}_j\mathscr{L}_{j+1}\mathscr{L}_j\mathbf{Y}
\cong\mathscr{L}_{j+1}\mathscr{L}_j\mathscr{L}_{j+1}\mathbf{Y}.
\end{gathered}
\end{equation*}

其结果是：我们得到群 $\ZZ^m\rtimes\mathrm{Br}_m$ 在长度为 $m$ 的%
\trExceptionalCollection{}的\trIsomorphism{}类之集合上的一个作用。子群 $\ZZ^m$
通过\trSuspension{}来作用；\trBraidGroup{} $\mathrm{Br}_m$ 的标准\trGenerator{}%
$\beta_1,\ldots,\beta_{m-1}$ 通过初等\trMutation{}来作用；而半直积则是相对于
对称群 $\mathrm{Sym}_m$ 的置换表示构作的，这里 $\mathrm{Sym}_m$ 被看作
$\mathrm{Br}_m$ 的一个\trQuotient。此外，若一个\trExceptionalCollection{}是
完整的，则在它的 $\ZZ^m\rtimes\mathrm{Br}_m$-\trTrajectory{}中的所有序列也都
如此。

这些证明都是围绕\trExactTriangle{}的简单论证，可以在\trTriangulated{}\trCategory{}%
的框架内完成而无需提及\trAinfty-结构；例如见 [66, Proposition 2.4.2]。我们要
强调，诸 $\mathscr{L}_k$ \emph{并非}\trFunctor：它们作用在 $\scr{A}$ 内由%
\trObject{}组成的序列上，而不作用在这些\trObject{}之间的\trMorphism{}上。
